\section{System Design}
\label{sec:design}

\textsc{Sentinel} is a set of cooperating services on one host, connected by
a Redis publish/subscribe bus (Fig.~\ref{fig:arch}). We describe the system as
built and evaluated, which differs in places from its original design; where
it does, we say so.

\begin{figure}[t]
\centering
\resizebox{\columnwidth}{!}{% SENTINEL architecture — the detection-to-response path as built.
% Solid boxes run on the appliance; dashed boxes are external services.
\begin{tikzpicture}[
  font=\footnotesize,
  stage/.style={draw, rounded corners=2pt, align=center, minimum height=9mm,
                minimum width=19mm, fill=blue!6},
  ext/.style={draw, dashed, rounded corners=2pt, align=center, minimum height=6.5mm,
              minimum width=19mm, fill=gray!8},
  bus/.style={-{Latex[length=1.8mm]}, thick},
  side/.style={-{Latex[length=1.5mm]}, densely dotted},
  lbl/.style={font=\scriptsize\ttfamily, inner sep=1.2pt}
]
% main path (y = 0), 34 mm pitch leaves room for the channel labels
\node[stage] (cap) at (0,0)    {Capture\\{\scriptsize Scapy}};
\node[stage] (ano) at (3.4,0)  {Anomaly\\{\scriptsize Kitsune, LSTM-AE}};
\node[stage] (cls) at (6.8,0)  {Classifier\\{\scriptsize Gemma 4 E2B}};
\node[stage] (rsp) at (10.2,0)  {Responder\\{\scriptsize policy graph}};
\draw[bus] (cap) -- node[lbl, above] {packets} (ano);
\draw[bus] (ano) -- node[lbl, above] {anomaly} (cls);
\draw[bus] (cls) -- node[lbl, above] {threat}  (rsp);

% enforcement row (y = -1.7)
\node[stage, fill=orange!12] (hp) at (0,-1.7)   {Honeypots\\{\scriptsize Cowrie / OpenCanary}};
\node[stage, fill=red!9]     (fw) at (10.2,-1.7) {nftables\\{\scriptsize block + guard}};
\draw[bus] (rsp) -- node[lbl, right] {actions} (fw);
\draw[bus] (hp) -- node[lbl, above] {attacker IP} (fw);

% local registry (on the appliance), above the main path
\node[stage, fill=green!7, minimum height=7mm] (reg) at (3.4,1.35) {Device registry};
\draw[side] (reg.east) -| node[lbl, above, pos=0.25] {vendor, CVEs} (cls.north);

% the appliance boundary
\node[draw, dotted, rounded corners=3pt, inner xsep=2.2mm, inner ysep=2.2mm,
      fit=(cap)(rsp)(hp)(fw)(reg)] (box) {};
\node[font=\scriptsize\itshape, anchor=north west] at (box.south west) {on the appliance};

% external services, above the boundary
\node[ext] (nvd)   at (3.4,2.75) {NIST NVD};
\node[ext] (intel) at (8.5,2.75) {OTX, AbuseIPDB};
\draw[side] (nvd) -- node[lbl, left, pos=0.3] {vendor name} (reg);
\draw[side] (intel.south) |- ++(0,-0.45) -| node[lbl, right, pos=0.75] {remote IP} (cls.north east);
\end{tikzpicture}
}
\caption{The detection-to-response path. Solid arrows are bus channels;
dashed boxes are the only components that contact external services.}
\label{fig:arch}
\end{figure}

\subsection{Capture and device context}
A capture service reads packets from the monitored interface with Scapy's
asynchronous sniffer and publishes a compact per-packet record (addresses,
ports, length, timestamp) on the bus. A discovery service builds a device
registry passively from the mDNS, SSDP, DHCP and ARP traffic it observes, maps
MAC prefixes to vendors, and queries NIST's National Vulnerability Database for
known CVEs by vendor. Active Nmap scanning exists but is triggered only on
demand from the dashboard.

\subsection{Anomaly detection}
For every packet the anomaly service computes the 115 Kitsune
features~\cite{mirsky2018kitsune}: damped statistics at five time scales
($\lambda \in \{5, 3, 1, 0.1, 0.01\}$) over five stream aggregations (source
MAC--IP pair, source IP, channel, channel jitter, and socket), including
two-dimensional statistics that relate each channel or socket to its reverse
direction. It keeps a sliding window of the last ten feature vectors per
device and scores each full window with a sequence autoencoder: a two-layer
LSTM~\cite{hochreiter1997lstm} encoder compresses the $10\times115$ window into
a 64-dimensional state, which is repeated across ten steps and decoded by a
second two-layer LSTM and a linear layer (153,651 parameters in total). The
anomaly score is the mean squared reconstruction error over the whole window,
computed with ONNX Runtime~\cite{onnxruntime} on min--max-scaled features, and
an event is raised when it exceeds a single threshold shared by all devices.

\subsection{LLM classification}
The classifier service enriches each anomaly event with the device's vendor
and known CVEs from the registry and with the reputation of the remote address
from AlienVault OTX~\cite{otx} and AbuseIPDB~\cite{abuseipdb}. It then asks
Gemma~4~E2B~\cite{gemma4}, served locally by Ollama~\cite{ollama} and quantised
to 4 bits, for a JSON object with an ATT\&CK tactic and technique identifier,
a confidence, a severity from 1 to 10, a one-sentence summary and a
recommended action. The model runs with Ollama's default sampling
(temperature~1.0), so its output for a given event is not deterministic.
Events whose response does not parse are dropped.

\subsection{Response and enforcement}
The response service is a two-node LangGraph~\cite{langgraph} state machine,
not a second language model: a policy node maps the LLM's severity to actions
through a declarative policy file, and a gate node holds the action for human
confirmation when the policy requires it. In the evaluated configuration,
severities 6--10 cause the event's indicator address---the destination of the
anomalous flow---to be blocked without confirmation, and lower severities are
logged. The policy also names a device-quarantine action for severities 9--10;
the responder currently enforces it as the same address block, and a
MAC-level quarantine module exists but is not invoked by the responder. Keeping the enforcement decision deterministic and
inspectable, rather than generated, is deliberate: the model proposes a
severity, and a policy the owner can read decides what happens.

Blocks are applied as nftables rules in a dedicated table, each tracked by its
rule handle so that it can be removed individually. A guard refuses to block
the appliance's own addresses, the default gateway, loopback, broadcast and
unspecified addresses, and address ranges, so that a misclassification cannot
cut the home off its own router. An XDP fast path is implemented but was not
enabled for this evaluation.

\subsection{Deception}
Cowrie~\cite{cowrie} (SSH and Telnet) and OpenCanary~\cite{opencanary} run in
containers. A honeypot service follows Cowrie's JSON log, records credentials
and commands from each session as a threat event, and sends the attacker's
address to the same guarded enforcement path.

\subsection{What leaves the appliance}
No packet payload, capture or feature vector is sent off the device, and the
LLM runs locally. Three lookups do leave it: the remote address of each
anomalous flow goes to OTX and AbuseIPDB, and device vendor names go to the
NVD. All three are optional, and a privacy-sensitive deployment can disable
them at some cost to context.

\subsection{Threat model}
In scope are network-observable attacks involving devices on the monitored
segment: compromised devices scanning or flooding (as in N-BaIoT), and inbound
credential attacks against exposed services (as seen by the honeypots). We
assume the attacker does not control the appliance and cannot tamper with its
view of the traffic. Out of scope are physical access, inspection of encrypted
payloads, supply-chain compromise, attacks on \textsc{Sentinel} itself, and
adaptive attackers who shape traffic to evade the detector.
